\subsection*{U.9 Sequential candidate admission on one paused Lift runtime}
\label{supp:sequential-lift-manager}

\paragraph{Question and preregistration.}
This bounded software experiment tests the manager loop directly: can it
withhold one candidate, preserve the actual pending boundary, check the
next candidate, and continue after the first passing decision?
Each of the three previously frozen Lift boundaries has one runtime object,
one mutable source-state object, and one executed old prefix.
The candidate priority is fixed before execution as released native,
saved normalized, then saved preference; selection stops at the first
complete strict-and-cooperative pass.
The separate packet is
\path{experiment/v64_lift_manager_loop_20261024/}.
Its \path{PLAN.json} fixes 70 exact input/code pins, the candidate order,
the three existing twelve-input workloads, and one scientific execution,
with a 120-second wall bound, 1\,GiB sampled RSS bound, 16\,MiB per-file
bound, and 128\,MiB output-tree bound.

An initial supervisor launch failed in a macOS address-space-limit
setter before the worker was executed. No checker or runtime episode
ran: the failure occurred in \code{preexec\_fn}, and the worker's output
directory was never created. The original plan, supervisor, and empty
streams are retained. A separately frozen \path{PREEXEC_AMENDMENT.json}
removes the unsupported address-space setter while retaining the RSS,
wall, and output limits, and places the actual execution in
\path{run02/}. The worker code is byte-identical across this amendment.
There is one scientific execution and no outcome-based retry.

\paragraph{Fresh checks and one-object selection.}
Before execution, both existing checker routes freshly validate the
complete old-mode closure: 50 nodes and 238 edges, with input coverage,
source safety, strictness, and cooperation all passing.
The runtime then executes the frozen two-edge prefix and pauses at the
actual pending state. At each candidate attempt, the manager rechecks
source and model bindings, saved-native/source predicate correspondence,
the exact ten source/monitor values, and the candidate's stored private
node and rank. Both routes traverse the entire candidate-rooted graph
and recompute coverage, safety, strictness, and cooperation.
Admission requires agreement and all five Boolean fields to be true.
No saved candidate verdict or certificate loader supplies the decision;
an audit hook rejects reads of the prior execution certificates.
Raw extraction-result files referenced by the inherited input manifest
are hashed for integrity only, not interpreted as admission verdicts.

All three native attempts are complete and strict but noncooperative
(154 nodes, 1,170 edges), so the manager withholds them. Each following
normalized attempt is complete, safe, strict, and cooperative
(154 nodes, 1,176 edges), so it passes and is installed:
\begin{center}
\begin{tabular}{lccc}
\toprule
Reached native node/rank & First decision & Installed node/rank & Commits\\
\midrule
524/1 & Withhold & 72/1 & 12\\
526/1 & Withhold & 74/1 & 12\\
756/1 & Withhold & 88/0 & 12\\
\bottomrule
\end{tabular}
\end{center}
Across each pair of checks, the runtime identity, source-state-object
identity, all source/monitor values, old-controller identity, private
memory, and commit count remain unchanged; the runtime stays paused.
The inherited installation adapter then checks the passing certificate's
exact graph digest, node, rank, and complete source/monitor match.
Only the selected controller object changes at installation.
The preference candidate is explicitly recorded as unattempted after
the first pass; this experiment makes no new claim about its verdict.

\paragraph{Continuation and scope.}
All three selected controllers execute their frozen twelve inputs.
The old floor-3 request is answered at post-update step~1 and its monitor
clears at step~2. No cancellation occurs, no genuinely pending obligation
remains, and the final executed transition is a repeatable self-loop
with all three monitors clear. The run has six fresh candidate checks,
three installations, 36 post-update commits, and 69 saved events; no
candidate check is unknown. The supervisor reports 0.650\,s overall
wall time and the worker reports 74,399,744 bytes peak RSS.
These are measurements of one instrumented run, not a latency benchmark.
\path{run02/output/RESULT.json} and \path{run02/output/EVENTS.jsonl} preserve every
attempt, before/after identity, certificate, installation, and trace.

This extends the separate candidate episodes of Section~U.4 with a
sequential rejection--admission loop on the same paused state.
It demonstrates enforcement by the software adapter using the existing
frozen graphs and predicates. It performs no synthesis, contract
weakening, winning-region recomputation, source/monitor reset, or
initializer reapplication at handoff. The complete-graph checks establish
the candidate properties; the twelve-step workloads demonstrate the
recorded continuation after those decisions.

An independent saved-artifact audit checks all 69 events, exact object
identities, rooted graph coverage, safety, and strongly connected components;
it also compares the source predicates on all 1,048,576 source-state pairs
using separate truth-table evaluation and native-node parsing.
The audit confirms the sequential decisions and continuations without
rerunning the manager or synthesis. Native exporter metadata and process
execution provenance remain trust boundaries; the audit record is
\path{independent_audit/AUDIT_RESULT.json} in this packet.
The continuation release supplies selected unchanged records
under \path{records/manager/} and the qualified \code{manager} saved-replay
stage. It rechecks graphs, source predicates and the 69-event sequence;
it does not execute the manager worker or revalidate the omitted original
source/runtime pins and worker call-site/read-path provenance.
